\documentclass[11pt]{article}

% ===============================
% Geometry and Layout
% ===============================
\usepackage[margin=1in]{geometry}
\usepackage{setspace}
\setstretch{1.15}

% ===============================
% Fonts and Encoding (LuaLaTeX)
% ===============================
\usepackage{fontspec}
\setmainfont{Latin Modern Roman}

% ===============================
% Mathematics and Symbols
% ===============================
\usepackage{amsmath, amssymb, amsthm}
\usepackage{mathtools}

% ===============================
% Hyperlinks
% ===============================
\usepackage{hyperref}
\hypersetup{
  colorlinks=true,
  linkcolor=black,
  urlcolor=black
}

% ===============================
% Theorem Environments
% ===============================
\theoremstyle{definition}
\newtheorem{definition}{Definition}
\theoremstyle{plain}
\newtheorem{proposition}{Proposition}
\newtheorem{lemma}{Lemma}

% ===============================
% Metadata
% ===============================
\title{The Geometry of Outrage:\\
A Field-Theoretic Analysis of the Attention Economy}
\author{Flyxion}
\date{\today}

\begin{document}
\maketitle

\begin{abstract}
This essay analyzes the contemporary attention economy as a constrained socio-technical system whose dominant dynamics arise from engagement-based optimization under severe temporal and visibility constraints. Rather than treating sensationalism, outrage, and radicalization as moral failures or cultural pathologies, the paper advances a structural account in which these phenomena emerge as stable outcomes of field-level dynamics. Attention is modeled as a finite allocable resource operationalized through engagement metrics, which function as proxies for compressive informational load. Recommendation systems are shown to induce path-dependent trajectories through content space, producing incremental but directional escalation even in the absence of malicious intent.

Drawing on systems theory, information theory, and variational principles, the essay formalizes the “rabbit hole” effect as gradient-following dynamics in a distorted informational field. Outrage is analyzed as a stable informational attractor favored by short-term optimization, while integrative, semantically rich content is systematically suppressed due to temporal scale mismatch. The analysis further demonstrates why moral and agent-centric explanations fail to account for recurrence, stability, and cross-platform convergence.

The Relational Scalar–Vector–Entropy Plenum (RSVP) framework is introduced as a formal interpretive lens that unifies scalar entropy density, vector flow, and constraint geometry within a single explanatory schema. Under this mapping, the attention economy is revealed as a concrete instantiation of a general class of systems in which proxy optimization produces entropy-efficient but semantically fragile outcomes. The essay concludes by reframing journalism, moderation, and governance as problems of constraint design rather than content correction, arguing that the preservation of meaning in high-throughput informational environments depends on the deliberate engineering of boundary conditions rather than on individual virtue alone.
\end{abstract}


\section{Introduction: The Crisis of the Attention Market}

Over the past two decades, public discourse has undergone a structural transformation. What was once organized around relatively stable institutions of mediation—newspapers, broadcasters, universities, and professional associations—has been reconfigured into a competitive market for attention. In this market, visibility is scarce, time is fragmented, and informational success is measured not by coherence or durability, but by immediate engagement. The consequences of this shift are now widely visible: polarization, sensationalism, epistemic instability, and the routine amplification of extremal content.

Contemporary analyses often describe these outcomes in moral or cultural terms. They are framed as failures of civic virtue, as evidence of declining trust, or as the product of irresponsible actors exploiting new technologies. While such explanations capture important dimensions of the problem, they struggle to account for its regularity. Similar pathologies recur across platforms, political contexts, and content genres, even when participants hold divergent values and intentions. Systems designed to maximize connection reliably produce fragmentation. Platforms optimized for participation repeatedly amplify outrage.

This essay advances a different claim. The dynamics of the attention economy are not accidental, nor are they primarily the result of moral failure. They are the predictable consequences of a particular structural configuration. When information systems are optimized for engagement under tight temporal constraints, they induce directional flows through informational space. Certain forms of content become energetically favored, while others are systematically suppressed. Outrage, sensationalism, and conspiratorial sense-making emerge not as deviations from the system’s logic, but as stable outcomes within it.

To make this claim precise, the essay adopts a field-theoretic perspective. Attention is treated not as a psychological state or a moral capacity, but as a finite allocable resource measured indirectly through behavioral proxies. Engagement metrics act as signals within an optimization regime, shaping the distribution and amplification of content. Recommendation systems induce path-dependent trajectories that guide users through informational environments in ways that are locally rational yet globally distorting.

The central thesis is therefore structural: the attention economy can be understood as a constrained informational field in which optimization for short-term engagement produces entropy descent. Under these conditions, escalation is not an anomaly but a gravitational effect. Any attempt to address the resulting crises—whether in journalism, politics, or artificial intelligence—must begin by understanding the geometry that produces them.

\section{Analytical Framework and Scope}

Before proceeding, it is necessary to fix the ontology and scope of the analysis. Without doing so, critiques of the attention economy are prone to category errors, oscillating between moral condemnation, psychological speculation, and conspiracy narratives about platform intent. This essay explicitly rejects such approaches in favor of a structural dynamics account.

The system under consideration is the contemporary attention economy understood as a distributed socio-technical system. It is composed of interacting agents and infrastructures, including users, content creators, platforms, and algorithmic ranking systems. These components operate under measurable constraints, most notably finite user time, competitive visibility, and monetization pressures tied to engagement metrics. The system evolves through feedback loops in which past performance signals shape future exposure.

The theory developed here is not a moral theory. It does not seek to classify content as virtuous or corrupt, nor does it assign blame to individual users or creators. It is also not a psychological model of belief formation. The analysis makes no claims about the internal cognitive states of participants, nor about the truth or falsity of the beliefs they may come to hold. Finally, it is not an ideological critique. The dynamics described apply across political orientations and content domains.

Instead, this essay offers a structural account of informational dynamics. Its explanatory targets are selection pressures, feedback mechanisms, and emergent patterns. The guiding question is not why particular actors behave as they do, but why certain behaviors and content forms systematically outcompete others under given constraints. Individual intent is treated as a local variable embedded within a global geometry that strongly conditions outcomes.

Several exclusions follow directly from this framing. The analysis does not assume coordinated malice on the part of platforms. It does not posit hidden conspiracies or deliberate attempts to radicalize users. Nor does it claim that outrage is inherently irrational or that participants are uniquely gullible. The goal is to explain recurrence, stability, and scale-independence. Any account that relies on contingent bad actors or cultural decline cannot meet this standard.

By fixing the system boundary in this way, the essay establishes the conditions under which its claims hold. What follows should be read not as a diagnosis of individual failures, but as an analysis of a field whose dynamics remain stable across changes in personnel, platform, and content.

\subsection*{Assumptions and Scope Conditions}

To make the analytical commitments explicit, we state the core assumptions underlying the framework.

\begin{definition}[System Boundary]
The attention economy is defined as the coupled system of content production, algorithmic selection, and user interaction operating under engagement-based optimization.
\end{definition}

\begin{definition}[Optimization Constraint]
An optimization constraint is a rule or objective function that determines which signals are amplified, prioritized, or suppressed within the system.
\end{definition}

\begin{lemma}[Independence from Intent]
If system-level outcomes recur across heterogeneous actors and platforms, then those outcomes cannot be fully explained by individual intent.
\end{lemma}

\begin{proof}
Assume outcomes depend primarily on intent. Then varying intent distributions should produce qualitatively different global dynamics. Empirically, similar escalation patterns recur across platforms and contexts with divergent actor intentions. Therefore intent alone is insufficient to explain the observed regularities.
\end{proof}

These assumptions delimit the explanatory scope of the essay and justify the transition to a field-theoretic description in subsequent sections.

\section{Core Concepts and Definitions}

A rigorous structural analysis requires that key terms be defined operationally rather than rhetorically. The concepts introduced here are not intended to capture subjective experience or moral evaluation. They are functional definitions suitable for describing system-level dynamics.

\begin{definition}[Attention]
Attention is a finite allocable resource corresponding to time and perceptual bandwidth expended by users within an informational system. It is not directly observable, but is approximated through behavioral proxies such as dwell time, click-through rates, and interaction frequency.
\end{definition}

\begin{definition}[Engagement Metric]
An engagement metric is an operational signal used by platforms to estimate attention allocation. Engagement metrics serve as optimization targets for algorithmic ranking and recommendation systems.
\end{definition}

\begin{definition}[Outrage]
Outrage is defined as a high-arousal, low-latency affective response that resolves rapidly into measurable engagement signals. It is characterized by intensity rather than duration.
\end{definition}

\begin{definition}[Algorithmic Selection]
Algorithmic selection is the automated ranking, filtering, and recommendation of content based on past performance signals, typically engagement metrics, under constraints of limited display capacity.
\end{definition}

\begin{definition}[Radicalization]
Radicalization, as used in this essay, refers to directional drift in content exposure over time toward increasingly extreme or polarized material. It does not denote belief conversion, ideological commitment, or psychological transformation.
\end{definition}

These definitions intentionally avoid normative language. They are designed to describe observable system behavior rather than internal states or ethical judgments.

\section{The Attention Economy as a Constrained Field}

With these definitions in place, the attention economy can be modeled as a constrained field. The term ``field'' is used here in a structural rather than physical sense. It denotes a state space in which content, affect, and visibility are distributed, and within which dynamics unfold under explicit optimization constraints.

The field consists of a high-dimensional space whose coordinates include, at minimum, affective intensity, temporal compression, semantic complexity, and amplification potential. Each piece of content occupies a region of this space, and user trajectories correspond to sequences of points generated through interaction with recommendation systems.

Crucially, this field is not neutral. It is shaped by constraints imposed by platform design. The dominant constraint is engagement-based optimization: content is ranked and propagated according to its ability to elicit measurable attention within short temporal windows. This constraint effectively reshapes the geometry of the space, amplifying some regions while suppressing others.

Regions associated with high affective intensity and rapid response become energetically favorable. Content that produces immediate engagement signals is preferentially sampled and redistributed. By contrast, regions associated with high semantic complexity but delayed payoff become dynamically inaccessible. They are not forbidden, but they are systematically underrepresented due to the mismatch between their informational structure and the system’s sampling resolution.

This produces a distortion analogous to a steepening of gradients within the field. Small differences in affective intensity translate into large differences in visibility. As optimization proceeds iteratively, the field evolves toward configurations in which high-arousal regions dominate the accessible landscape.

It is important to emphasize that this distortion does not require explicit design choices favoring sensationalism. It arises as an emergent property of the optimization target. When engagement metrics are treated as proxies for value, the system implicitly privileges compressive informational forms over integrative ones. This bias propagates across scales, influencing creator behavior, user experience, and institutional viability.

The field-theoretic framing therefore shifts explanatory focus. Rather than asking why certain content succeeds, it asks which regions of the field are rendered stable under given constraints. Rather than attributing outcomes to individual preference, it identifies selection pressures that operate independently of intent.

\subsection*{Field Structure and Constraints}

We formalize the preceding discussion in abstract terms.

Let $\mathcal{C}$ denote the space of content representations, and let $\mathcal{E}: \mathcal{C} \rightarrow \mathbb{R}^{+}$ be an engagement functional mapping content to expected engagement.

\begin{definition}[Constraint Functional]
A constraint functional is an objective function $\mathcal{E}$ such that content propagation probability is monotonic in $\mathcal{E}$.
\end{definition}

\begin{proposition}[Constraint-Induced Distortion]
If $\mathcal{E}$ is sensitive primarily to short-term affective intensity, then content regions with delayed semantic payoff will be under-sampled regardless of their long-term informational value.
\end{proposition}

\begin{proof}
Sampling occurs at discrete time intervals optimized for rapid signal detection. Content whose engagement signal integrates slowly will produce lower instantaneous $\mathcal{E}$ values at sampling times. Therefore, such content will be systematically deprioritized under monotonic optimization.
\end{proof}

This result does not depend on the specific form of $\mathcal{E}$, only on its temporal resolution. The distortion is therefore structural rather than contingent.

\section{Incremental Drift and Directional Escalation}

A central empirical observation in contemporary analyses of the attention economy is that escalation rarely occurs through sudden discontinuities. Users are not typically propelled from mainstream content into extremism in a single step, nor do creators abruptly abandon restraint in favor of sensationalism. Instead, the dominant pattern is incremental. Each transition appears locally reasonable, while the aggregate trajectory becomes increasingly distorted. This section formalizes that observation as a property of path-dependent systems operating under engagement-based optimization.

Recommendation systems do not present content in isolation. They generate sequences. Each recommendation is conditioned on prior interactions, and each subsequent interaction updates the system’s estimate of what content is likely to maximize engagement. From the perspective of the user, this appears as a series of choices. From the perspective of the system, it is a process of iterative local optimization.

The crucial feature of such systems is that local optimality does not imply global neutrality. When transitions are selected according to immediate engagement gain, the resulting paths acquire directionality. Small biases compound. Each step is justified by marginal improvement, yet the sequence as a whole drifts toward regions of the content space that maximize short-term affective density.

This phenomenon can be described as incremental drift. At each iteration, the system selects a content item that lies slightly closer to a local engagement maximum than the previous one. The magnitude of each step may be negligible in isolation. However, because the system repeatedly applies the same selection rule, the cumulative effect becomes substantial. Direction emerges not from any single choice, but from the repeated application of a biased rule.

Path dependence follows directly. Once a user has traversed a sequence of increasingly high-arousal content, the system’s future recommendations are conditioned on that history. Alternative trajectories become less visible, not because they are forbidden, but because they are less compatible with the system’s learned engagement profile. The informational cost of reversal increases with path length.

This introduces hysteresis into the system. Returning to earlier regions of the content space requires more than simply reversing the most recent step. It requires undoing accumulated conditioning effects in the recommendation model. In practical terms, this means that exposure trajectories are not easily reversible. The system remembers where attention has been, and that memory shapes where it can go next.

Importantly, none of these dynamics require ideological alignment between platform, creator, and user. Incremental escalation is compatible with heterogeneous intentions. Creators may believe they are merely refining presentation. Users may believe they are exploring adjacent topics. Platforms may believe they are neutral optimizers. Yet the global effect is directional because the selection pressure is asymmetric.

This structural account explains why escalation appears gradual rather than abrupt. Sudden transitions would be easily detected and resisted. Incremental drift, by contrast, exploits continuity. Each step lies within the tolerance of local rationality. The system therefore produces outcomes that feel unchosen, even though no single step is coercive.

Understanding escalation as incremental drift has important implications. It shifts the explanatory burden away from moments of conversion or persuasion and toward the geometry of permissible paths. It also clarifies why interventions focused on isolated content items are ineffective. Removing a single node does not alter the curvature of the space through which paths are drawn.

Incremental drift is thus not a failure of awareness or intent. It is a predictable outcome of systems that optimize locally under globally distorting constraints.

\subsection*{Path Dependence and Hysteresis}

We formalize incremental drift using abstract dynamical systems language.

Let $\mathcal{C}$ be the content space and let $x_t \in \mathcal{C}$ denote the content exposed at time $t$. Let the recommendation rule be given by
\[
x_{t+1} = \arg\max_{y \in \mathcal{N}(x_t)} \mathcal{E}(y),
\]
where $\mathcal{N}(x_t)$ is a neighborhood of candidate content and $\mathcal{E}$ is the engagement functional.

\begin{proposition}[Incremental Drift]
If $\mathcal{E}$ has a nonzero gradient in a given region of $\mathcal{C}$, then repeated local maximization induces directional drift along that gradient.
\end{proposition}

\begin{proof}
At each step, the selection rule chooses a point with higher $\mathcal{E}$ value within $\mathcal{N}(x_t)$. If $\nabla \mathcal{E} \neq 0$, then the expected displacement has a component in the direction of $\nabla \mathcal{E}$. Iteration yields cumulative displacement.
\end{proof}

\begin{lemma}[Hysteresis]
If the recommendation rule conditions on past exposure history, then reversing a trajectory requires more than reversing the most recent step.
\end{lemma}

\begin{proof}
Conditioning introduces state dependence on the entire history $\{x_0, \dots, x_t\}$. The transition kernel at time $t+1$ therefore differs from that at earlier times, even for identical $x_t$, preventing simple reversibility.
\end{proof}

These properties are generic to locally optimized sequential systems and do not depend on the semantic content of $\mathcal{C}$.

\section{Outrage as a Stable Informational Attractor}

The persistence and ubiquity of outrage within the attention economy is often interpreted as a cultural pathology or a deviation from rational discourse. Such interpretations obscure a more fundamental property of the system. Outrage is not merely present; it is stable. Across platforms, topics, and political orientations, high-arousal content reliably re-emerges as a dominant form. This regularity suggests that outrage functions as an informational attractor within the constrained field described earlier.

In this essay, the term ``attractor'' is used in a precise, system-theoretic sense. An attractor is a region of a state space toward which trajectories converge under the system’s dynamics. Importantly, an attractor need not be globally optimal, normatively desirable, or epistemically accurate. It need only be dynamically stable: once trajectories enter its basin, they tend to remain there unless acted upon by external forces.

Outrage qualifies as such an attractor because it satisfies the stability conditions imposed by engagement-based optimization. High-arousal content generates rapid, high-amplitude engagement signals. These signals are immediately legible to the system’s sampling mechanisms and are therefore preferentially amplified. As a result, trajectories that approach outrage-rich regions of the content space experience reinforcement rather than correction.

This reinforcement operates at multiple levels. At the level of users, repeated exposure to high-arousal content conditions future recommendations, narrowing the range of subsequent material. At the level of creators, content that elicits outrage is rewarded with visibility, incentivizing replication and refinement. At the level of platforms, aggregated engagement validates the optimization target, reinforcing the underlying selection rule. Each layer contributes to the stabilization of outrage as a dominant mode.

Crucially, the stability of outrage does not depend on topic or ideology. The same affective structure can be instantiated through political commentary, cultural conflict, health misinformation, or entertainment media. What matters is not the propositional content but the affective profile. This explains why attempts to suppress specific narratives often fail. Removing one instantiation of outrage does not eliminate the attractor; it merely displaces trajectories toward alternative instantiations.

The stability of outrage also explains its resistance to correction. Introducing factual information into a region dominated by high-arousal dynamics does not alter the underlying gradient. On the contrary, corrective interventions may themselves be absorbed into the attractor if they generate comparable affective intensity. The system does not distinguish between outrage directed at an external target and outrage directed at outrage itself. Both register as high-density signals.

This observation clarifies why outrage is often mistaken for preference. Users may report fatigue, dissatisfaction, or distress, yet continue to engage with outrage-inducing content. From a structural perspective, this is not paradoxical. Engagement does not measure endorsement or enjoyment; it measures responsiveness. Outrage persists because it reliably produces responses under tight temporal constraints.

Understanding outrage as a stable informational attractor reframes the challenge it poses. The problem is not that outrage exists, nor that individuals are drawn to it in isolation. The problem is that the system’s geometry renders outrage dynamically privileged. As long as engagement-based optimization defines the field, outrage will continue to function as a sink for attention, regardless of countervailing intentions.

Any attempt to address the dominance of outrage must therefore contend with its stability. Moral condemnation, content removal, and user education may alter local conditions, but they do not change the shape of the field. Without structural modification, trajectories will continue to converge toward the same regions. Stability, not deviance, is the defining feature.

\subsection*{Attractors and Stability}

We formalize the notion of an informational attractor in abstract terms.

Let $\mathcal{C}$ denote the content space and let the dynamics be given by a transition operator $T: \mathcal{C} \rightarrow \mathcal{C}$ induced by recommendation and engagement feedback.

\begin{definition}[Attractor]
A subset $A \subset \mathcal{C}$ is an attractor if there exists a neighborhood $U$ of $A$ such that for all $x \in U$, the sequence $\{T^n(x)\}$ converges to $A$ as $n \rightarrow \infty$.
\end{definition}

\begin{proposition}[Outrage Stability]
If content within a region $A$ of $\mathcal{C}$ produces engagement signals that exceed those of neighboring regions under the same sampling constraints, then $A$ constitutes an attractor under engagement-based optimization.
\end{proposition}

\begin{proof}
By assumption, $\mathcal{E}(x) > \mathcal{E}(y)$ for $x \in A$ and $y \in U \setminus A$. The transition operator $T$ therefore preferentially maps states in $U$ into $A$. Iteration yields convergence.
\end{proof}

\begin{lemma}[Topic-Independence]
If engagement depends primarily on affective intensity rather than propositional content, then multiple disjoint regions of $\mathcal{C}$ may instantiate equivalent attractors.
\end{lemma}

\begin{proof}
Affective intensity defines equivalence classes in $\mathcal{C}$ under $\mathcal{E}$. Any region sharing the same affective profile will satisfy the attractor condition.
\end{proof}

These results demonstrate that outrage is dynamically stable under the stated constraints, independent of content domain or intent.

\section{Why Moral and Cultural Critiques Are Insufficient}

Public discourse surrounding the attention economy frequently adopts a moral register. Outrage is framed as a failure of character, sensationalism as a lapse in journalistic ethics, and radicalization as a consequence of intellectual irresponsibility or cultural decay. These critiques are not without merit. Individuals do make choices, institutions do set norms, and values do matter. However, as explanatory frameworks, moral and cultural accounts are insufficient. They do not explain why the same patterns recur across platforms, populations, and ideological contexts, nor why well-intentioned actors consistently produce outcomes they themselves disavow.

The central limitation of moral critique lies in its focus on agency at the wrong level of description. Moral explanations operate at the level of individual intention and responsibility. They ask why people choose poorly, why institutions fail to uphold standards, or why audiences reward harmful content. Yet the phenomena under consideration persist even when individual intentions vary widely. Journalists committed to accuracy, users seeking information, and platforms claiming neutrality all participate in systems that nonetheless converge on sensationalism and escalation.

Cultural explanations encounter a related problem. Accounts that attribute outrage to polarization, tribalism, or declining trust identify important contextual factors, but they remain historically contingent. They struggle to explain why similar dynamics emerge in distinct cultural settings or why newly introduced platforms reproduce familiar pathologies with remarkable speed. If culture were the primary driver, one would expect greater variation in outcomes across contexts. Instead, one observes structural convergence.

The insufficiency of these explanations becomes especially clear when considering predictability. Moral critiques can condemn outcomes after the fact, but they do not generate testable predictions about system behavior under changing conditions. They cannot, for example, explain why altering interface timing or recommendation cadence produces measurable changes in escalation rates, while appeals to civility do not. Structural accounts, by contrast, predict such sensitivity because they identify the mechanisms through which constraints shape trajectories.

This does not imply that moral responsibility is irrelevant. Rather, it implies that moral responsibility is constrained by geometry. Individuals act within environments that make certain behaviors easy and others costly. When a system reliably rewards affective compression and penalizes deliberation, moral exhortation functions as a weak force opposing a strong gradient. Expecting virtue alone to counteract such gradients is akin to expecting frictionless motion uphill.

Recognizing the limits of moral explanation also clarifies why debates about intent are often unproductive. Whether platforms intend to amplify outrage or whether creators consciously exploit outrage is secondary to the fact that outrage amplification occurs regardless. Intent may modulate the rate of descent, but it does not determine the direction. Structural incentives dominate.

The analytical consequence of this observation is decisive. To explain recurrence, stability, and scale, one must shift from agent-centered narratives to field-level causation. Moral language may still play a role in motivating reform, but it cannot serve as the foundation of explanation. Without a structural account, reform efforts remain reactive and fragile, addressing symptoms rather than causes.

\subsection*{Agent-Level vs. Field-Level Causation}

We distinguish two levels of causation relevant to the analysis.

\begin{definition}[Agent-Level Causation]
Agent-level causation refers to outcomes attributable to the intentions, beliefs, or decisions of individual actors.
\end{definition}

\begin{definition}[Field-Level Causation]
Field-level causation refers to outcomes that arise from the structure of the system, independent of the specific intentions of individual actors.
\end{definition}

\begin{proposition}[Recurrence Criterion]
If a pattern recurs across systems with heterogeneous agents and intentions, then its primary cause must be field-level rather than agent-level.
\end{proposition}

\begin{proof}
Assume agent-level causation dominates. Then varying agent distributions should produce qualitatively different outcomes. Empirical recurrence across platforms and contexts contradicts this assumption. Therefore, field-level causation is necessary to explain recurrence.
\end{proof}

\begin{lemma}[Moral Insufficiency]
Moral exhortation alone cannot reliably alter outcomes determined by field-level constraints.
\end{lemma}

\begin{proof}
Moral exhortation operates by modifying agent-level preferences. If field-level constraints impose strong gradients, small changes in preference are insufficient to reverse system trajectories.
\end{proof}

These results justify the transition from moral critique to structural design in the subsequent section.

\section{Implications for Institutional Design}

Once the dynamics of the attention economy are understood as field-level phenomena, the implications for institutional design become unavoidable. Institutions that mediate public meaning—news organizations, platforms, regulatory bodies, and civic infrastructures—do not merely transmit content. They shape the geometry within which informational trajectories unfold. Their design choices function as boundary conditions that determine which paths are stable, which are costly, and which are effectively inaccessible.

From this perspective, journalism is not primarily a content practice but a structural one. Editorial standards such as verification, contextualization, proportionality, and correction operate as constraints on how information may propagate. These practices slow down amplification, introduce friction into feedback loops, and preserve semantic depth over time. In an engagement-optimized environment, such constraints are not neutral. They impose energetic costs that place institutions at a competitive disadvantage relative to actors unconstrained by similar commitments.

This observation clarifies why journalistic restraint increasingly appears anomalous within the attention market. It is not that restraint has lost its value, but that the surrounding field penalizes it. When platforms optimize for rapid affective response, institutions that privilege delayed understanding must expend additional resources to maintain visibility. The cost of integrity is therefore structural rather than moral. It reflects the work required to oppose prevailing gradients.

Transparency initiatives, often framed as trust-building exercises, acquire a different significance in this context. Making editorial processes visible does not merely signal good faith. It externalizes the labor involved in maintaining semantic architecture. By exposing how stories are selected, checked, and revised, institutions reveal the scaffolding that resists collapse into sensationalism. Transparency, in this sense, is not performance. It is load-bearing structure.

Similar considerations apply to moderation and governance. Policies that focus exclusively on removing harmful content after it appears address endpoints rather than trajectories. They intervene once escalation has already occurred. By contrast, design choices that affect pacing, recommendation depth, and amplification thresholds operate upstream, altering the curvature of the field itself. These choices shape not only what content is seen, but how users move through informational space.

The institutional challenge, then, is not to eliminate outrage or prevent disagreement. Conflict is not a pathology. The challenge is to prevent runaway concentration of attention into extremal attractors that suppress alternative modes of sense-making. Institutions capable of sustaining plural trajectories—multiple coexisting paths through content space—exhibit resilience. Those that permit unbounded acceleration toward high-arousal regions do not.

Importantly, this reframing avoids the false dichotomy between regulation and free expression. Constraint design is not synonymous with censorship. Every system already imposes constraints; the question is whether they are implicit and unmanaged or explicit and deliberate. Designing for semantic resilience requires acknowledging this fact rather than denying it.

Institutional design, understood in this way, becomes a problem of ethics in the literal sense: the shaping of conditions under which collective life unfolds. Values are not appended to systems from the outside. They are instantiated through constraint geometry. Whether meaning persists or collapses depends less on declarations of principle than on the structures built to support them.

\subsection*{Constraint Engineering}

We formalize the role of institutions as constraint designers.

Let $\mathcal{C}$ denote the content space and let $T$ be the transition operator induced by platform design.

\begin{definition}[Boundary Condition]
A boundary condition is a design parameter that restricts or shapes the set of allowable transitions in $\mathcal{C}$.
\end{definition}

\begin{proposition}[Path Sensitivity]
Altering boundary conditions that govern transition structure has a greater effect on long-term system behavior than altering isolated content items.
\end{proposition}

\begin{proof}
Long-term behavior is determined by repeated application of $T$. Modifying individual content items affects only isolated states, whereas modifying boundary conditions alters $T$ itself, changing the evolution of all trajectories.
\end{proof}

\begin{lemma}[Constraint Cost]
If a boundary condition reduces engagement velocity, then institutions enforcing it must supply compensatory resources to remain viable.
\end{lemma}

\begin{proof}
Engagement-based optimization penalizes slower transitions. Institutions that impose friction incur opportunity costs relative to unconstrained actors, requiring external support to persist.
\end{proof}

These results formalize the intuition that integrity-preserving institutions perform counter-incentive labor within an engagement-optimized field.

\section{RSVP Interpretation: Attention as a Scalar--Vector--Entropy Field}

The preceding analysis has proceeded without invoking any particular formal framework beyond general systems and field-theoretic language. This section introduces the Relational Scalar--Vector--Entropy Plenum (RSVP) framework as an interpretive formalization of the dynamics already described. The purpose of this move is not to revise the argument, but to make explicit the structural correspondences that have implicitly guided it.

RSVP is a general field-theoretic framework designed to describe systems in which scalar quantities, directional flows, and boundary constraints jointly determine irreversible dynamics. It is agnostic to substrate. While originally developed in physical and informational contexts, its formal structure applies equally to socio-technical systems whose behavior is governed by optimization under constraint. The attention economy analyzed in this essay constitutes one such system.

Within the RSVP framework, the dominant variables of the attention economy map cleanly onto the three core components of the theory. Engagement metrics correspond to a scalar entropy field. They encode compressive informational load: how much affective or perceptual impact is concentrated into a given temporal interval. High-arousal content registers as regions of high entropy density, not because it is disordered, but because it produces rapid, high-amplitude signals under constrained sampling.

Recommendation systems correspond to vector flows induced by gradients in this scalar field. Users are not modeled as choosing among static options, but as traversing trajectories through content space. Each recommendation step represents a small displacement along an entropy gradient. Over time, these displacements integrate into directional drift. Radicalization, in this formulation, is not a property of content nodes, but of paths.

Algorithmic incentive regimes correspond to constraint geometry. Visibility thresholds, monetization rules, and engagement decay rates define the boundary conditions under which scalar and vector dynamics unfold. Algorithms are not treated as agents with intent, but as operators that carve the space of admissible trajectories. Creators and users adapt to this geometry regardless of ideology or motivation.

This mapping clarifies several results established earlier. The stability of outrage follows from its status as a high-entropy attractor. Incremental escalation follows from gradient-following dynamics under local optimization. The failure of moral exhortation follows from the dominance of field-level constraints over agent-level intent. Editorial restraint appears as counter-entropic labor because it requires sustained energy to oppose prevailing gradients.

Importantly, introducing RSVP at this stage does not add new empirical claims. It adds explanatory compression. By unifying scalar density, directional flow, and constraint geometry within a single formal schema, RSVP explains why the same pathologies recur across platforms and contexts, and why reform efforts that ignore structure remain ineffective.

The attention economy, understood through this lens, is not uniquely pathological. It is a particular instantiation of a more general phenomenon: systems optimized for proxy signals under irreversible constraints tend to generate entropy-efficient but semantically fragile attractors. Media systems merely render this dynamic visible at scale.

This interpretation also clarifies the broader relevance of the analysis. The same structural logic appears in governance systems optimized for short-term approval, in financial markets optimized for liquidity, and in artificial intelligence systems optimized for reward proxies. In each case, scalar optimization induces vector drift under constraint, producing outcomes that diverge from intended values. RSVP provides a language for recognizing these failures as homologous rather than accidental.

\subsection*{RSVP Mapping}

We summarize the formal correspondences introduced above.

\begin{definition}[Scalar Entropy Field]
A scalar entropy field is a mapping $\rho: \mathcal{C} \rightarrow \mathbb{R}^{+}$ encoding compressive informational load under constrained sampling.
\end{definition}

\begin{definition}[Vector Flow]
A vector flow is a rule assigning directional displacement $\vec{V}(x) \propto \nabla \rho(x)$ governing trajectory evolution in $\mathcal{C}$.
\end{definition}

\begin{definition}[Constraint Geometry]
Constraint geometry is the set of boundary conditions that restrict admissible trajectories and shape the effective topology of $\mathcal{C}$.
\end{definition}

\begin{proposition}[RSVP Consistency]
The dynamics of the attention economy described in this essay are isomorphic to an RSVP system with engagement-defined entropy, recommendation-induced flow, and algorithmic boundary conditions.
\end{proposition}

\begin{proof}
Engagement metrics define $\rho$. Recommendation systems implement $\vec{V}$. Platform incentives define admissible transitions. All prior results follow from these identifications.
\end{proof}

This formalization does not depend on the semantic content of $\mathcal{C}$ and therefore applies across domains.

\section{Conclusion}

This essay has analyzed the contemporary attention economy as a distributed socio-technical system governed by engagement-based optimization under severe temporal and visibility constraints. Attention was treated not as a moral faculty or psychological disposition, but as a finite allocable resource operationalized through behavioral proxies. Recommendation systems were modeled as sequential selection mechanisms that generate trajectories through content space rather than isolated choices. Platforms, creators, and users were situated within a shared geometry shaped by measurable incentives rather than by intent.

The central constraint identified throughout the analysis is the use of short-term engagement metrics as optimization targets. This constraint privileges compressive informational forms that produce rapid affective responses while systematically disadvantaging integrative forms whose semantic value accumulates over longer timescales. Under repeated application, this optimization reshapes the informational field, steepening gradients toward high-arousal regions and suppressing access to alternative trajectories.

The emergent outcome of this configuration is directional escalation. Incremental drift replaces discrete conversion, outrage stabilizes as a dominant attractor, and reversal becomes increasingly costly due to path dependence and hysteresis. These dynamics recur across platforms and contexts because they arise from field-level structure rather than from contingent moral or cultural failures. Individual intent modulates local behavior but does not determine global outcomes.

Institutional responses that operate solely at the level of content or individual responsibility are therefore insufficient. Durable change requires modification of boundary conditions that govern trajectory formation. Journalism, moderation, and governance function structurally as forms of constraint engineering. Their viability depends on whether truth-preserving paths can be maintained against prevailing gradients without being energetically suppressed.

The RSVP framework was introduced not as an alternative explanation, but as a formal consolidation of these results. By naming scalar entropy, vector flow, and constraint geometry explicitly, it clarifies why similar failures emerge wherever proxy optimization governs irreversible systems. The attention economy is one instance of a broader class of dynamics in which optimization for convenient signals produces coherent but fragile outcomes.

The implication is not that outrage must be eliminated, nor that conflict must be suppressed. It is that systems which fail to encode semantic resilience into their structure will reliably converge on entropy-efficient attractors regardless of the values of their participants. Once this geometry is recognized, neutrality is no longer a defensible stance. Constraint design becomes the medium through which collective meaning either persists or collapses.

\section*{Appendices}

\appendix
\section{Appendix A: RSVP Field Equations}

\subsection{Prerequisites}

Let $(M, g)$ be a smooth $n$-dimensional manifold.
Let $t \in \mathbb{R}$ denote a global time parameter.
Assume sufficient regularity for all fields.

\subsection{Fundamental Fields}

Define the RSVP fields:

\[
\Phi : M \times \mathbb{R} \to \mathbb{R}
\qquad \text{(scalar density field)}
\]

\[
\vec{v} : M \times \mathbb{R} \to TM
\qquad \text{(vector flow field)}
\]

\[
S : M \times \mathbb{R} \to \mathbb{R}_{\ge 0}
\qquad \text{(entropy density field)}
\]

\subsection{Derived Quantities}

Material derivative:
\[
D_t := \partial_t + \vec{v} \cdot \nabla
\]

Entropy flux:
\[
\vec{J}_S := S \vec{v}
\]

Effective potential:
\[
V(\Phi, S) := \alpha \Phi^2 + \beta S \Phi + \gamma S^2
\]

\subsection{Field Equations (Eulerian Form)}

Scalar density evolution:
\[
D_t \Phi = - \lambda \nabla \cdot \vec{v}
\]

Vector flow evolution:
\[
D_t \vec{v} = - \nabla \Phi - \nabla S + \nu \Delta \vec{v}
\]

Entropy balance:
\[
\partial_t S + \nabla \cdot \vec{J}_S = \sigma(\Phi, \vec{v}) \ge 0
\]

with entropy production:
\[
\sigma := \eta \|\nabla \vec{v}\|^2 + \kappa (\nabla \Phi \cdot \nabla S)
\]

\subsection{Action Functional}

Define the RSVP action:
\[
\mathcal{A}[\Phi, \vec{v}, S] :=
\int dt \int_M
\left(
\frac{1}{2} \|\vec{v}\|^2
- V(\Phi,S)
- \chi S \log S
\right)
\, d\mu_g
\]

\subsection{Variational Derivatives}

Euler--Lagrange equations:
\[
\frac{\delta \mathcal{A}}{\delta \Phi} = 0,
\qquad
\frac{\delta \mathcal{A}}{\delta \vec{v}} = 0,
\qquad
\frac{\delta \mathcal{A}}{\delta S} = 0
\]

Entropy constraint:
\[
\delta S \ge 0
\]

\subsection{Canonical Momenta}

\[
\pi_\Phi := \frac{\partial \mathcal{L}}{\partial (\partial_t \Phi)} = 0
\]

\[
\vec{\pi}_v := \frac{\partial \mathcal{L}}{\partial (\partial_t \vec{v})} = \vec{v}
\]

\[
\pi_S := \frac{\partial \mathcal{L}}{\partial (\partial_t S)} = 0
\]

Primary constraints:
\[
\pi_\Phi \approx 0,
\qquad
\pi_S \approx 0
\]

\subsection{Hamiltonian Density}

\[
\mathcal{H} =
\frac{1}{2} \|\vec{\pi}_v\|^2
+ V(\Phi,S)
+ \chi S \log S
\]

Total Hamiltonian:
\[
H_T = \int_M \mathcal{H}\, d\mu_g
+ \lambda_\Phi \pi_\Phi
+ \lambda_S \pi_S
\]

\subsection{Poisson Structure}

Canonical brackets:
\[
\{\Phi(x), \pi_\Phi(y)\} = \delta(x-y)
\]

\[
\{v_i(x), \pi_{v_j}(y)\} = \delta_{ij}\delta(x-y)
\]

\[
\{S(x), \pi_S(y)\} = \delta(x-y)
\]

\subsection{Irreversibility Condition}

Entropy monotonicity:
\[
\frac{d}{dt} \int_M S \, d\mu_g \ge 0
\]

Time-reversal symmetry broken iff:
\[
\sigma > 0
\]

\section{Appendix B: Variational Constraints and Irreversible Histories}

\subsection{Configuration Space}

Let
\[
\mathcal{Q} := \{\Phi, \vec{v}, S\}
\]
denote the configuration space of RSVP fields.

Let
\[
\mathcal{H} := \{ \gamma : \mathbb{R} \to \mathcal{Q} \}
\]
denote the space of histories.

\subsection{Constraint Manifold}

Define admissible histories:
\[
\mathcal{H}_{\mathrm{adm}} :=
\left\{
\gamma \in \mathcal{H} \;\middle|\;
\partial_t S \ge 0
\right\}
\]

Define constraint functional:
\[
\mathcal{C}[\gamma] :=
\int_M \sigma(\Phi,\vec{v},S)\, d\mu_g
\]

Admissibility condition:
\[
\mathcal{C}[\gamma] \ge 0
\]

\subsection{Constrained Action Principle}

Define constrained action:
\[
\mathcal{A}_{\mathrm{adm}} :=
\inf_{\gamma \in \mathcal{H}_{\mathrm{adm}}}
\mathcal{A}[\gamma]
\]

Variations restricted to:
\[
\delta \gamma \in T_\gamma \mathcal{H}_{\mathrm{adm}}
\]

\subsection{Projection Operator}

Define projection:
\[
\mathcal{P} : \mathcal{H} \to \mathcal{H}_{\mathrm{adm}}
\]

such that:
\[
\mathcal{A}[\mathcal{P}(\gamma)]
\le
\mathcal{A}[\gamma]
\]

\subsection{Coarse-Graining Operator}

Define coarse-graining:
\[
\mathcal{G}_\Lambda :
\mathcal{H} \to \mathcal{H}_\Lambda
\]

with resolution parameter $\Lambda$.

Non-invertibility:
\[
\mathcal{G}_\Lambda^{-1} \;\; \text{undefined}
\]

\subsection{Effective Action}

Define effective action:
\[
\mathcal{A}_\Lambda :=
\mathcal{A} \circ \mathcal{G}_\Lambda
\]

Monotonicity:
\[
\mathcal{A}_{\Lambda_1} \le \mathcal{A}_{\Lambda_2}
\quad
\text{for}
\quad
\Lambda_1 < \Lambda_2
\]

\subsection{History Selection Functional}

Define selection weight:
\[
W[\gamma] :=
\exp\left(
- \int_0^T \sigma(\gamma(t)) \, dt
\right)
\]

Support:
\[
\mathrm{supp}(W) \subset \mathcal{H}_{\mathrm{adm}}
\]

\subsection{Irreversibility}

Define time reversal:
\[
\mathcal{T} : \gamma(t) \mapsto \gamma(-t)
\]

Violation condition:
\[
\mathcal{T}(\mathcal{H}_{\mathrm{adm}})
\cap
\mathcal{H}_{\mathrm{adm}}
=
\varnothing
\]

\subsection{Admissible Path Integral}

Define path integral:
\[
Z :=
\int_{\mathcal{H}_{\mathrm{adm}}}
\mathcal{D}\gamma
\;
e^{-\mathcal{A}[\gamma]}
\]

Normalization:
\[
\int_{\mathcal{H}_{\mathrm{adm}}} W[\gamma] = 1
\]

\subsection{Loss of Detailed Balance}

Detailed balance violated iff:
\[
W[\gamma] \neq W[\mathcal{T}\gamma]
\]

Entropy arrow:
\[
\frac{d}{dt} S_{\mathrm{tot}} > 0
\]

\section{Appendix C: Categorical Structure of RSVP}

\subsection{Base Category}

Define category $\mathcal{R}$ such that:

Objects:
\[
\mathrm{Ob}(\mathcal{R}) :=
\{ \gamma \in \mathcal{H}_{\mathrm{adm}} \}
\]

Morphisms:
\[
\mathrm{Hom}_{\mathcal{R}}(\gamma_1,\gamma_2) :=
\left\{
f : \gamma_1 \to \gamma_2 \;\middle|\;
f \text{ preserves admissibility}
\right\}
\]

Identity:
\[
\mathrm{id}_\gamma : \gamma \to \gamma
\]

Composition:
\[
g \circ f : \gamma_1 \to \gamma_3
\]

Associativity:
\[
(h \circ g) \circ f = h \circ (g \circ f)
\]

\subsection{Admissibility-Preserving Morphisms}

Constraint preservation:
\[
f^\ast(\sigma_{\gamma_2}) \le \sigma_{\gamma_1}
\]

Entropy monotonicity:
\[
S(\gamma_2) \ge S(\gamma_1)
\]

\subsection{Symmetric Monoidal Structure}

Tensor product:
\[
\otimes : \mathcal{R} \times \mathcal{R} \to \mathcal{R}
\]

Unit object:
\[
\mathbb{I} := \text{vacuum history}
\]

Associator:
\[
\alpha_{\gamma_1,\gamma_2,\gamma_3}
:
(\gamma_1 \otimes \gamma_2)\otimes \gamma_3
\to
\gamma_1 \otimes (\gamma_2 \otimes \gamma_3)
\]

Braiding:
\[
\beta_{\gamma_1,\gamma_2}
:
\gamma_1 \otimes \gamma_2
\to
\gamma_2 \otimes \gamma_1
\]

Symmetry:
\[
\beta_{\gamma_2,\gamma_1} \circ \beta_{\gamma_1,\gamma_2}
=
\mathrm{id}
\]

\subsection{Functorial Coarse-Graining}

Define functor:
\[
\mathcal{G}_\Lambda :
\mathcal{R} \to \mathcal{R}_\Lambda
\]

Functoriality:
\[
\mathcal{G}_\Lambda(\mathrm{id}_\gamma)
=
\mathrm{id}_{\mathcal{G}_\Lambda(\gamma)}
\]

\[
\mathcal{G}_\Lambda(g \circ f)
=
\mathcal{G}_\Lambda(g)
\circ
\mathcal{G}_\Lambda(f)
\]

\subsection{Natural Transformations}

Resolution morphism:
\[
\eta_{\Lambda_1 \to \Lambda_2}
:
\mathcal{G}_{\Lambda_1}
\Rightarrow
\mathcal{G}_{\Lambda_2}
\quad
(\Lambda_1 > \Lambda_2)
\]

Commutativity:
\[
\eta_{\Lambda_2 \to \Lambda_3}
\circ
\eta_{\Lambda_1 \to \Lambda_2}
=
\eta_{\Lambda_1 \to \Lambda_3}
\]

\subsection{Limits and Colimits}

Projective limit:
\[
\gamma =
\varprojlim_{\Lambda}
\mathcal{G}_\Lambda(\gamma)
\]

Non-existence of inverse limit:
\[
\varinjlim_{\Lambda}
\mathcal{G}_\Lambda^{-1}
\quad
\text{undefined}
\]

\subsection{Entropy as Functorial Measure}

Define entropy functor:
\[
\mathcal{S} : \mathcal{R} \to (\mathbb{R}_{\ge 0}, \le)
\]

Monotonicity:
\[
f : \gamma_1 \to \gamma_2
\implies
\mathcal{S}(\gamma_1) \le \mathcal{S}(\gamma_2)
\]

\subsection{Gauge Equivalence}

Define equivalence:
\[
\gamma_1 \sim \gamma_2
\iff
\exists f,g :
\gamma_1 \rightleftarrows \gamma_2
\]

Gauge groupoid:
\[
\mathcal{Gpd}(\mathcal{R})
\]

\section{Appendix D: Sheaf-Theoretic Structure of RSVP}

\subsection{Base Space}

Let $(M,\tau)$ be a topological space.
Let $\mathcal{O}(M)$ denote the category of open sets of $M$.

\subsection{Presheaf of Histories}

Define presheaf:
\[
\mathcal{F} : \mathcal{O}(M)^{op} \to \mathbf{Set}
\]

Sections:
\[
\mathcal{F}(U) :=
\{ \gamma_U : U \times \mathbb{R} \to \mathcal{Q} \}
\]

Restriction maps:
\[
\rho_{UV} : \mathcal{F}(U) \to \mathcal{F}(V)
\quad
(V \subset U)
\]

Functoriality:
\[
\rho_{UU} = \mathrm{id}
\quad
\text{and}
\quad
\rho_{VW} \circ \rho_{UV} = \rho_{UW}
\]

\subsection{Sheaf Condition}

Cover:
\[
U = \bigcup_i U_i
\]

Matching:
\[
\rho_{U_i U_{ij}}(\gamma_i)
=
\rho_{U_j U_{ij}}(\gamma_j)
\]

Gluing:
\[
\exists!\;
\gamma \in \mathcal{F}(U)
\quad
\text{s.t.}
\quad
\rho_{U U_i}(\gamma)=\gamma_i
\]

\subsection{Admissible Subsheaf}

Define subsheaf:
\[
\mathcal{F}_{\mathrm{adm}} \subset \mathcal{F}
\]

Condition:
\[
\partial_t S_{\gamma_U} \ge 0
\quad
\forall \gamma_U \in \mathcal{F}_{\mathrm{adm}}(U)
\]

\subsection{Sheaf of Constraints}

Define constraint sheaf:
\[
\mathcal{C}(U) :=
\{ \sigma_U : U \to \mathbb{R}_{\ge 0} \}
\]

Compatibility:
\[
\rho_{UV}(\sigma_U) = \sigma_V
\]

\subsection{Obstruction Cocycle}

Čech cocycle:
\[
\omega_{ijk} \in Z^2(\mathcal{U},\mathcal{C})
\]

Obstruction condition:
\[
\omega_{ijk} \neq 0
\;\Rightarrow\;
\text{no global admissible section}
\]

\subsection{Derived Sections}

Derived global sections:
\[
\mathbb{R}\Gamma(M,\mathcal{F}_{\mathrm{adm}})
\]

Nontrivial cohomology:
\[
H^k(M,\mathcal{F}_{\mathrm{adm}}) \neq 0
\quad (k>0)
\]

\subsection{Topological Entropy Invariant}

Define entropy class:
\[
[S] \in H^0(M,\mathbb{R}_{\ge 0})
\]

Monotonicity:
\[
\frac{d}{dt}[S] \ge 0
\]

\subsection{Sheaf Morphisms}

Define morphism:
\[
\varphi : \mathcal{F}_{\mathrm{adm}} \to \mathcal{F}_{\mathrm{adm}}
\]

Commutativity:
\[
\rho_{UV} \circ \varphi_U
=
\varphi_V \circ \rho_{UV}
\]

\subsection{Stack Structure}

Define stack:
\[
\mathfrak{H} := [\mathcal{F}_{\mathrm{adm}}/\sim]
\]

Descent data:
\[
\gamma_i \sim \gamma_j \text{ on } U_{ij}
\]

Effective descent:
\[
\gamma \in \mathfrak{H}(U)
\]

\subsection{Entropy Obstruction}

Obstruction class:
\[
[\omega_S] \in H^1(M,\mathcal{C})
\]

Vanishing condition:
\[
[\omega_S] = 0
\iff
\exists \text{ global entropy-coherent history}
\]

\section{Appendix E: Extensions to CLIO, TARTAN, and Semantic Drift}

\subsection{CLIO Dynamics}

Define CLIO state:
\[
\Xi := (\gamma, \mathcal{G}_\Lambda, \mathcal{P})
\]

Define CLIO operator:
\[
\mathcal{C}\mathcal{L} : \Xi_t \mapsto \Xi_{t+1}
\]

Constraint:
\[
\mathcal{C}\mathcal{L}(\Xi_t) \in \mathcal{H}_{\mathrm{adm}}
\]

Fixed point:
\[
\Xi^\ast = \mathcal{C}\mathcal{L}(\Xi^\ast)
\]

Stability:
\[
\| \Xi_{t+1} - \Xi^\ast \| \le \kappa \| \Xi_t - \Xi^\ast \|,
\quad 0 < \kappa < 1
\]

\subsection{TARTAN Structure}

Define tiling cover:
\[
\mathcal{U} := \{ U_i \subset M \}
\]

Define trajectory-aware section:
\[
\gamma_i := \gamma|_{U_i \times [t_i,t_{i+1}]}
\]

Define recursive tiling operator:
\[
\mathcal{T} : \mathcal{F}_{\mathrm{adm}}(U_i) \to
\mathcal{F}_{\mathrm{adm}}(\{U_{i_k}\})
\]

Recursion depth:
\[
\mathcal{T}^n(\gamma_i)
\]

\subsection{Annotated Noise Field}

Define noise field:
\[
\eta : M \times \mathbb{R} \to \mathbb{R}
\]

Define annotation map:
\[
\mathcal{A} : \eta \mapsto \Sigma
\quad
(\Sigma \in \mathcal{S}_{\mathrm{meta}})
\]

Annotated perturbation:
\[
\tilde{\gamma} := \gamma + \epsilon \eta_\Sigma
\]

Constraint:
\[
\partial_t S_{\tilde{\gamma}} \ge 0
\]

\subsection{TARTAN Gluing Condition}

Compatibility:
\[
\rho_{U_i U_{ij}}(\gamma_i)
=
\rho_{U_j U_{ij}}(\gamma_j)
\]

Annotated consistency:
\[
\mathcal{A}(\eta_i)|_{U_{ij}}
=
\mathcal{A}(\eta_j)|_{U_{ij}}
\]

Tiled colimit:
\[
\Gamma_{\mathrm{TARTAN}}
:=
\operatorname{colim}_{i}
(\gamma_i + \eta_i)
\]

\subsection{Trajectory Preservation}

Define trajectory functional:
\[
\mathcal{T}_\gamma := \{\gamma(t)\}_{t \in \mathbb{R}}
\]

Preservation condition:
\[
\mathcal{T}_{\Gamma_{\mathrm{TARTAN}}}
\sim
\mathcal{T}_\gamma
\]

Entropy bound:
\[
\int_M S_{\Gamma_{\mathrm{TARTAN}}}
\le
\int_M S_{\gamma} + \delta
\]

\subsection{Conflict-Free Recursive Merge}

Define merge operator:
\[
\mu : (\gamma_i + \eta_i) \otimes (\gamma_j + \eta_j)
\to
\gamma_{ij}
\]

Associativity:
\[
\mu(\mu(x,y),z)=\mu(x,\mu(y,z))
\]

Commutativity:
\[
\mu(x,y)=\mu(y,x)
\]

Idempotence:
\[
\mu(x,x)=x
\]

\subsection{Semantic Metric Space}

Define semantic probability measures:
\[
\mu_t \in \mathcal{P}(\mathcal{C})
\]

Define Wasserstein metric:
\[
d_W(\mu_1,\mu_2)
=
\inf_{\pi \in \Pi(\mu_1,\mu_2)}
\int_{\mathcal{C}\times\mathcal{C}}
\|x-y\|\, d\pi(x,y)
\]

\subsection{Semantic Drift}

Define drift functional:
\[
\Delta_S(\gamma_t,\gamma_{t+1})
=
d_W(\mu_t,\mu_{t+1})
\]

Monotonicity:
\[
\Delta_S \ge 0
\]

Bounded drift:
\[
\Delta_S \le \varepsilon
\]

\subsection{Alignment Fixed Points}

Define alignment operator:
\[
\mathcal{A}_{\mathrm{align}}
:
\mathcal{H}_{\mathrm{adm}} \to \mathcal{H}_{\mathrm{adm}}
\]

Fixed point:
\[
\mathcal{A}_{\mathrm{align}}(\gamma)=\gamma
\]

Terminal condition:
\[
\partial_t S_\gamma = 0
\Rightarrow
\gamma = \gamma_\infty
\]

\subsection{Proxy Collapse Criterion}

Define proxy functional:
\[
R : \mathcal{H} \to \mathbb{R}
\]

Collapse condition:
\[
\nabla R \parallel \nabla S
\]

Degeneracy:
\[
\exists \gamma_1 \neq \gamma_2 :
R(\gamma_1)=R(\gamma_2),
\quad
S(\gamma_1) \neq S(\gamma_2)
\]

\subsection{Entropy-Respecting Computation}

Define computation functor:
\[
\mathcal{E} : \mathcal{R} \to \mathcal{R}
\]

Entropy monotonicity:
\[
\mathcal{S}(\mathcal{E}(\gamma)) \ge \mathcal{S}(\gamma)
\]

Admissibility:
\[
\mathcal{E}(\mathcal{H}_{\mathrm{adm}})
\subseteq
\mathcal{H}_{\mathrm{adm}}
\]

\section{Appendix F: Spherepop Mapping}

\subsection{Spherepop State Space}

Define Spherepop universe:
\[
\mathcal{B} := \{ b_i \}_{i \in I}
\]

Define bubble state:
\[
b_i := (\phi_i, \vec{u}_i, s_i)
\]

State space:
\[
\mathcal{B} \subset \mathbb{R} \times \mathbb{R}^n \times \mathbb{R}_{\ge 0}
\]

\subsection{RSVP--Spherepop Correspondence}

Define projection:
\[
\Pi : (\Phi, \vec{v}, S) \mapsto \{ b_i \}
\]

Componentwise:
\[
\phi_i := \int_{U_i} \Phi \, d\mu
\]

\[
\vec{u}_i := \frac{1}{|U_i|}\int_{U_i} \vec{v} \, d\mu
\]

\[
s_i := \int_{U_i} S \, d\mu
\]

\subsection{Bubble Adjacency Graph}

Define adjacency:
\[
G := (\mathcal{B}, E)
\]

Edge condition:
\[
(b_i,b_j) \in E
\iff
U_i \cap U_j \neq \varnothing
\]

\subsection{Local Update Operator}

Define Spherepop update:
\[
\mathcal{S}_t : b_i^t \mapsto b_i^{t+1}
\]

Scalar update:
\[
\phi_i^{t+1} = \phi_i^t - \lambda \sum_{j \sim i} (\vec{u}_i - \vec{u}_j)
\]

Vector update:
\[
\vec{u}_i^{t+1} =
\vec{u}_i^t
- \nabla_i \phi_i
- \nabla_i s_i
\]

Entropy update:
\[
s_i^{t+1} = s_i^t + \sigma_i
\quad
\sigma_i \ge 0
\]

\subsection{Discrete Entropy Production}

Define local production:
\[
\sigma_i :=
\eta \sum_{j \sim i} \|\vec{u}_i - \vec{u}_j\|^2
\]

Monotonicity:
\[
s_i^{t+1} \ge s_i^t
\]

\subsection{Bubble Merge Operator}

Define merge:
\[
\mu_b : b_i \otimes b_j \to b_{ij}
\]

Associativity:
\[
\mu_b(\mu_b(b_i,b_j),b_k)
=
\mu_b(b_i,\mu_b(b_j,b_k))
\]

Commutativity:
\[
\mu_b(b_i,b_j)=\mu_b(b_j,b_i)
\]

Idempotence:
\[
\mu_b(b_i,b_i)=b_i
\]

\subsection{Bubble Split Operator}

Define split:
\[
\Delta_b : b_i \to \{b_{i_1},b_{i_2}\}
\]

Conservation:
\[
\phi_{i_1}+\phi_{i_2}=\phi_i
\]

\[
s_{i_1}+s_{i_2} \ge s_i
\]

\subsection{Spherepop Category}

Define category:
\[
\mathcal{S}\mathrm{ph}
\]

Objects:
\[
\mathrm{Ob}(\mathcal{S}\mathrm{ph}) := \mathcal{B}
\]

Morphisms:
\[
\mathrm{Hom}(b_i,b_j) :=
\{ \mu_b, \Delta_b, \mathcal{S}_t \}
\]

\subsection{Functor to RSVP}

Define functor:
\[
\mathcal{F}_{\mathrm{SP}} : \mathcal{S}\mathrm{ph} \to \mathcal{R}
\]

Preservation:
\[
\mathcal{S}(\mathcal{F}_{\mathrm{SP}}(b_i))
\ge
\mathcal{S}(b_i)
\]

\subsection{Discrete Trajectories}

Define bubble trajectory:
\[
\mathcal{T}_{b_i} := \{ b_i^t \}_{t \in \mathbb{N}}
\]

Lift:
\[
\mathcal{T}_{b_i} \mapsto \gamma \in \mathcal{H}_{\mathrm{adm}}
\]

\subsection{Spherepop Alignment}

Define alignment condition:
\[
\mathcal{S}_t(b_i)=b_i
\quad \forall i
\]

Global equilibrium:
\[
\sum_i s_i = \max
\]

\subsection{Spherepop Gauge Equivalence}

Define equivalence:
\[
b_i \sim b_j
\iff
\phi_i=\phi_j,\;
\vec{u}_i=\vec{u}_j,\;
s_i=s_j
\]

Quotient:
\[
\mathcal{B}/\sim
\]

\subsection{Terminal Configuration}

Define terminal state:
\[
\mathcal{B}_\infty :=
\{ b_i : \vec{u}_i = 0 \}
\]

Entropy maximum:
\[
\sum_i s_i = \sup
\]

\section{Appendix G: MAGI --- Manifold-Aligned Generative Inference}

\subsection{Latent Manifold}

Let $(\mathcal{Z}, g_{\mathcal{Z}})$ be a smooth Riemannian manifold.

Embedding:
\[
\iota : \mathcal{Z} \hookrightarrow \mathbb{R}^N
\]

Probability density:
\[
p(z) \in \mathcal{P}(\mathcal{Z})
\]

\subsection{Observation Space}

Let $\mathcal{X}$ be observation space.

Generative map:
\[
G : \mathcal{Z} \to \mathcal{X}
\]

Inference map:
\[
I : \mathcal{X} \to \mathcal{Z}
\]

Consistency:
\[
I \circ G \approx \mathrm{id}_{\mathcal{Z}}
\]

\subsection{Alignment Constraint}

Define alignment functional:
\[
\mathcal{L}_{\mathrm{align}} :=
\int_{\mathcal{Z}}
\| \nabla_{\mathcal{Z}} \log p(z) \|_{g_{\mathcal{Z}}}^2
\, d\mu_{\mathcal{Z}}
\]

MAGI constraint:
\[
\nabla_{\mathcal{Z}} \log p(z)
\in
T_z \mathcal{Z}
\]

\subsection{Geodesic Prediction}

Define geodesic flow:
\[
\frac{D}{dt}\dot{z}(t)=0
\]

Prediction operator:
\[
\hat{z}_{t+1} :=
\exp_{z_t}(\dot{z}_t)
\]

\subsection{MAGI Objective}

Define total objective:
\[
\mathcal{J}_{\mathrm{MAGI}} :=
\mathbb{E}_{x \sim \mathcal{X}}
\left[
\| x - G(I(x)) \|^2
+
\lambda \mathcal{L}_{\mathrm{align}}
\right]
\]

\subsection{Entropy Coupling (RSVP)}

Define pullback entropy:
\[
S(z) := - \log p(z)
\]

Alignment condition:
\[
\nabla S(z) \perp \ker(dG_z)
\]

\subsection{Vector Flow (Inference)}

Define inference flow:
\[
\vec{V}_{\mathrm{MAGI}} :=
- \nabla_{g_{\mathcal{Z}}} S(z)
\]

Constraint:
\[
\vec{V}_{\mathrm{MAGI}} \in T_z \mathcal{Z}
\]

\subsection{Trajectory Consistency}

Define trajectory:
\[
\mathcal{T}_z := \{ z(t) \}_{t \ge 0}
\]

Consistency:
\[
z(t+\Delta t)
=
\exp_{z(t)}(\Delta t \cdot \vec{V}_{\mathrm{MAGI}})
\]

\subsection{Semantic Stability}

Define semantic metric:
\[
d_{\mathcal{Z}}(z_1,z_2)
\]

Bound:
\[
d_{\mathcal{Z}}(z_t,z_{t+1}) \le \epsilon
\]

\subsection{MAGI--CLIO Coupling}

Define coupled state:
\[
\Omega := (\Xi, z)
\]

Update:
\[
(\Xi_{t+1},z_{t+1})
=
(\mathcal{C}\mathcal{L}(\Xi_t),
\exp_{z_t}(\vec{V}_{\mathrm{MAGI}}))
\]

Admissibility:
\[
(\Xi_{t+1},z_{t+1}) \in \mathcal{H}_{\mathrm{adm}} \times \mathcal{Z}
\]

\subsection{MAGI--TARTAN Interaction}

Define tile-local latent:
\[
z_i := z|_{U_i}
\]

Annotated noise:
\[
z_i \mapsto z_i + \eta_i
\]

Constraint:
\[
\nabla S(z_i+\eta_i) \in T_{z_i}\mathcal{Z}
\]

\subsection{MAGI--Spherepop Projection}

Define discretization:
\[
\pi_{\mathrm{MAGI}} : \mathcal{Z} \to \mathcal{B}
\]

Componentwise:
\[
\phi_i := \mathbb{E}_{z_i}[\Phi(z_i)]
\]

\[
\vec{u}_i := \mathbb{E}_{z_i}[\nabla G(z_i)]
\]

\[
s_i := \mathbb{E}_{z_i}[S(z_i)]
\]

\subsection{Alignment Fixed Points}

Define MAGI fixed point:
\[
\nabla S(z^\ast)=0
\]

Global alignment:
\[
(z^\ast,\gamma^\ast)
\quad
\text{s.t.}
\quad
\partial_t S_{\gamma^\ast}=0
\]

\subsection{Failure Mode (Off-Manifold Drift)}

Drift condition:
\[
\nabla S(z) \notin T_z \mathcal{Z}
\]

Projection error:
\[
\| z - \Pi_{\mathcal{Z}}(z) \| > \delta
\]

\subsection{Terminal Condition}

Terminal latent:
\[
z_\infty :=
\arg\min_{z \in \mathcal{Z}} S(z)
\]

Consistency:
\[
G(z_\infty)=x_\infty
\]

Entropy minimum:
\[
S(z_\infty)=\inf
\]

\begin{thebibliography}{99}

\bibitem{Simon1971}
Simon, H. A. (1971).
Designing organizations for an information-rich world.
In M. Greenberger (Ed.), \textit{Computers, communications, and the public interest} (pp. 37--72).
Baltimore: Johns Hopkins University Press.

\bibitem{Shannon1948}
Shannon, C. E. (1948).
A mathematical theory of communication.
\textit{Bell System Technical Journal}, 27(3), 379--423.

\bibitem{CoverThomas2006}
Cover, T. M., \& Thomas, J. A. (2006).
\textit{Elements of Information Theory}.
Hoboken, NJ: Wiley-Interscience.

\bibitem{Tufekci2015}
Tufekci, Z. (2015).
Algorithmic harms beyond Facebook and Google: Emergent challenges of computational agency.
\textit{Colorado Technology Law Journal}, 13, 203--218.

\bibitem{Tufekci2018}
Tufekci, Z. (2018).
YouTube, the great radicalizer.
\textit{The New York Times}, March 10.

\bibitem{Pariser2011}
Pariser, E. (2011).
\textit{The Filter Bubble: What the Internet Is Hiding from You}.
New York: Penguin Press.

\bibitem{Wu2016}
Wu, T. (2016).
\textit{The Attention Merchants: The Epic Scramble to Get Inside Our Heads}.
New York: Knopf.

\bibitem{Zuboff2019}
Zuboff, S. (2019).
\textit{The Age of Surveillance Capitalism}.
New York: PublicAffairs.

\bibitem{Anderson1972}
Anderson, P. W. (1972).
More is different.
\textit{Science}, 177(4047), 393--396.

\bibitem{Arthur1994}
Arthur, W. B. (1994).
\textit{Increasing Returns and Path Dependence in the Economy}.
Ann Arbor: University of Michigan Press.

\bibitem{Scheffer2009}
Scheffer, M. (2009).
\textit{Critical Transitions in Nature and Society}.
Princeton: Princeton University Press.

\bibitem{Bak1996}
Bak, P. (1996).
\textit{How Nature Works: The Science of Self-Organized Criticality}.
New York: Copernicus.

\bibitem{Lanczos1970}
Lanczos, C. (1970).
\textit{The Variational Principles of Mechanics}.
Toronto: University of Toronto Press.

\bibitem{PrigogineStengers1984}
Prigogine, I., \& Stengers, I. (1984).
\textit{Order Out of Chaos: Man's New Dialogue with Nature}.
New York: Bantam Books.

\bibitem{Barabasi2002}
Barabási, A.-L. (2002).
\textit{Linked: The New Science of Networks}.
Cambridge, MA: Perseus Publishing.

\bibitem{Villani2003}
Villani, C. (2003).
\textit{Topics in Optimal Transportation}.
Providence, RI: American Mathematical Society.

\bibitem{Goodhart1975}
Goodhart, C. A. E. (1975).
Problems of monetary management: The U.K. experience.
In \textit{Papers in Monetary Economics}.
Sydney: Reserve Bank of Australia.

\bibitem{Strathern1997}
Strathern, M. (1997).
Improving ratings: Audit in the British university system.
\textit{European Review}, 5(3), 305--321.

\bibitem{Benkler2018}
Benkler, Y., Faris, R., \& Roberts, H. (2018).
\textit{Network Propaganda: Manipulation, Disinformation, and Radicalization in American Politics}.
Oxford: Oxford University Press.

\bibitem{Goldberg2025}
Goldberg, M. (2025).
The conspiracy theory that’s consuming the right.
\textit{The New York Times}.

\bibitem{LaFrance2025}
LaFrance, A. (2025).
The prophecies of Q; The Epstein files.
\textit{The Atlantic}.

\bibitem{Berenson1992}
Berenson, E. (1992).
\textit{The Trial of Madame Caillaux: Scandal and Politics in the Belle Époque}.
Ithaca, NY: Cornell University Press.

\bibitem{Proust1913}
Proust, M. (1913--1927).
\textit{À la recherche du temps perdu}.
Paris: Grasset.

\bibitem{WarzelPetersen2021}
Warzel, C., \& Petersen, A. H. (2021).
\textit{Out of Office: Unlocking the Power and Potential of Hybrid Work}.
New York: Knopf.

\bibitem{FrumWarzel2025Video}
Frum, D., \& Warzel, C. (2025, December 31).
Facts vs. clicks: How algorithms reward extremism.
\textit{The David Frum Show x Galaxy Brain}.
The Atlantic.

\bibitem{Amodei2016}
Amodei, D., Olah, C., Steinhardt, J., Christiano, P., Schulman, J., \& Mané, D. (2016).
Concrete problems in AI safety.
arXiv:1606.06565.

\bibitem{Russell2019}
Russell, S. (2019).
\textit{Human Compatible: Artificial Intelligence and the Problem of Control}.
New York: Viking.

\bibitem{Landauer1961}
Landauer, R. (1961).
Irreversibility and heat generation in the computing process.
\textit{IBM Journal of Research and Development}, 5(3), 183--191.

\bibitem{Jaynes2003}
Jaynes, E. T. (2003).
\textit{Probability Theory: The Logic of Science}.
Cambridge: Cambridge University Press.

\bibitem{Friston2010}
Friston, K. (2010).
The free-energy principle: A unified brain theory?
\textit{Nature Reviews Neuroscience}, 11(2), 127--138.

\bibitem{CowanCRDT2023}
Cowan, M., et al. (2023).
Conflict-free replicated data types and distributed integrity.
\textit{Proceedings of the ACM}.

\end{thebibliography}

\end{document}

